\section{Anwendung und Benutzeroberfläche}

\subsection{Architektur der Streamlit-Anwendung}
\begin{itemize}
  \item \textbf{Framework:} Streamlit für schnelle interaktive Web-Anwendung ohne Frontend-Framework-Komplexität. % Repo-Evidenz prüfen
  \item \textbf{Navigation:} Sidebar-basiertes Page-Routing. Nutzer wählen Seiten über Selectbox. % Repo-Evidenz prüfen
  \item \textbf{Session State:} Streamlit \texttt{st.session\_state} hält Filter und Einstellungen über Neuladen. % Repo-Evidenz prüfen
  \item \textbf{Bootstrap:} \texttt{src/app/} orchestriert App-Startup, Auto-Import und Database-Resolver. % Repo-Evidenz prüfen
\end{itemize}

\subsection{Hauptseiten}
\begin{itemize}
  \item \textbf{Dashboard:} % Repo-Evidenz prüfen
    \begin{itemize}
      \item KPI-Reihen: Anzahl Trades, Gate-Passer, durchschnittlicher Score, Marktkapitalisierung-Aggregation.
      \item Charts:
        \begin{itemize}
          \item Horizontale Balkendiagramme für Buy/Sell nach Sektor.
          \item Trend-Chart für Netto-Signale über Zeit.
          \item Top-Companies und Top-Insiders nach Aktivität.
        \end{itemize}
      \item Filter: Datumsbereich (30 Tage per Default). % Repo-Evidenz prüfen
    \end{itemize}
  \item \textbf{Trade-Explorer (Hauptarbeitsflaeche):} % Repo-Evidenz prüfen
    \begin{itemize}
      \item Operative Filterfläche:
        \begin{itemize}
          \item Symbol (LIKE-Suche).
          \item Insider-Name.
          \item Richtung (Buy/Sell/Alle).
          \item Gate-Status (PASS/PENDING/FAIL/Alle).
          \item Validierungs-Status (VALID/INVALID/PRICE\_INVALID/Alle).
          \item Mindest-Score (Schieber, 0--100).
          \item Datumsbereich.
          \item Akkumulierungs-Modus: Gruppierung oder Einzeltrades.
        \end{itemize}
      \item Tabelle mit Paginierung (50/100/200 Reihen pro Seite). % Repo-Evidenz prüfen
        \begin{itemize}
          \item Spalten: Symbol, Filing Date, Reporting Name, Qty, Price, Trade Value, Gate Status, Score, Score Class.
          \item In Advanced Mode: zusätzliche Spalten (CIK, Market Cap, Industry, Validation Status).
        \end{itemize}
      \item Drilldown: Klick auf Trade öffnet Detailansicht. % Repo-Evidenz prüfen
    \end{itemize}
  \item \textbf{Unternehmen:} % Repo-Evidenz prüfen
    \begin{itemize}
      \item Unternehmens-Übersicht: Name, Ticker, Market Cap, Industrie, Trade Count, letzter Trade. % Repo-Evidenz prüfen
      \item Drilldown in Unternehmens-Detail-Seite mit Trade-Historie und Aggregationen. % Repo-Evidenz prüfen
    \end{itemize}
  \item \textbf{Methodik:} % Repo-Evidenz prüfen
    \begin{itemize}
      \item Dokumentation der Pipeline (Graphviz-Diagramm). % Repo-Evidenz prüfen
      \item Raw/Clean-Trennung erklärt.
      \item Validierung und Gate-Regeln textlich beschrieben.
      \item Score-Komponenten und Gewichtungen. % Interne Projektspezifikation
      \item Transaction-Code-Klassifikation.
    \end{itemize}
  \item \textbf{Einstellungen:} % Repo-Evidenz prüfen
    \begin{itemize}
      \item Konfigurierbare Gate-Schwellen (Mindestwert, Form Type).
      \item Score-Schwellen (PASS-Grenze, HOLD-Grenze).
      \item Import-Intervalle (für Auto-Import).
      \item Einstellungen werden in MongoDB persistiert. % Repo-Evidenz prüfen
    \end{itemize}
  \item \textbf{Admin:} % Repo-Evidenz prüfen
    \begin{itemize}
      \item Import-Steuerung: manueller Import mit Status-Reporting. % Repo-Evidenz prüfen
      \item Scheduler-Verwaltung: Auto-Import ein/aus, Tests.
      \item Datenbank-Status: Raw-Store-Zähler (MongoDB), Clean-Store-Info (MySQL).
      \item Öffentliche Freigabe (Public Share via Cloudflared).
      \item Raw~Clean-Sync (optionaler Re-Sync aus MongoDB ohne neue API-Calls).
    \end{itemize}
\end{itemize}

\subsection{UI-Komponenten und Interaktivität}
\begin{itemize}
  \item \textbf{Filterung:}
    \begin{itemize}
      \item Text-Eingaben (Symbol, Name) mit Wildcard-Matching.
      \item Selectboxes für Klassifikations-Felder (Gate Status, Richtung).
      \item Datumsbereichs-Picker (date\_input mit Start/End).
      \item Zahlenschieber für Score-Schwellen und Mindestwert.
      \item Einklappbares Filter-Panel zur UI-Platzoptimierung. % Repo-Evidenz prüfen
    \end{itemize}
  \item \textbf{Tabellen:} % Repo-Evidenz prüfen
    \begin{itemize}
      \item Sortierbar nach allen Spalten (standardmäßig Filing Date descending). % Repo-Evidenz prüfen
      \item Zentrierte Ausrichtung für tabellarische Daten (häufig in Analysetools). % Repo-Evidenz prüfen
      \item Highlighting nach Status (PASS grün, FAIL rot, PENDING gelb). % Repo-Evidenz prüfen
      \item Row-Selection für Drill-Down. % Repo-Evidenz prüfen
    \end{itemize}
  \item \textbf{Charts:}
    \begin{itemize}
      \item Balkendiagramme (matplotlib, seaborn oder Streamlit native).
      \item Zeitreihen-Liniendiagramme für Trends.
      \item Farben nach Kategorie (Sektor, Status). % Repo-Evidenz prüfen
    \end{itemize}
  \item \textbf{KPI-Rendering:}
    \begin{itemize}
      \item \texttt{render\_kpi\_row()} zeigt Metriken in Spalten-Layout. % Repo-Evidenz prüfen
      \item Werte können formatiert sein (Währung, Prozent, Zahl). % Repo-Evidenz prüfen
    \end{itemize}
\end{itemize}

\subsection{Statusdarstellung und Barrierefreiheit}
\begin{itemize}
  \item \textbf{Textliche Darstellung von Status:} % Repo-Evidenz prüfen
    \begin{itemize}
      \item Nicht ausschließlich Farbe: Jeder Status-Badge trägt auch Label-Text (\texttt{PASS}, \texttt{FAIL}, \texttt{PENDING}). % Repo-Evidenz prüfen
      \item WCAG-Konformität: Farbe allein ist nicht das einzige Erkennungsmittel. % Externe Quelle nötig
    \end{itemize}
  \item \textbf{Kontrast und Lesbarkeit:}
    \begin{itemize}
      \item Schriftgroessen und Farbpaletten orientieren sich an gaengigen Accessibility-Richtlinien. % TODO: pruefen
      \item Text-Alternativen für Charts und Icons.
    \end{itemize}
  \item \textbf{Navigation und Struktur:}
    \begin{itemize}
      \item Klare Seitentitel und Struktur.
      \item Sidebar Navigation ist konsistent.
      \item Tastatur-Navigation wird unterstuetzt. % TODO: pruefen
    \end{itemize}
\end{itemize}

\subsection{Advanced Mode}
\begin{itemize}
  \item \textbf{Funktionalität:} Toggle in Sidebar für zusätzliche Details. % Repo-Evidenz prüfen
  \item \textbf{Effekt:}
    \begin{itemize}
      \item Erweiterte Tabellenspalten: CIK, Profile Status, Technical Metrics.
      \item Detaillierte Scoring-Komponenten sichtbar.
      \item Fehler- und Warnungsdetails werden angezeigt.
    \end{itemize}
\end{itemize}
